\documentclass[11pt]{article}
\usepackage{icmlko}
\begin{document}

\title{두 번째로 깊어진다 --- 되돌린 근거가 그 일을 못 재는 자였다}
\author{월드모델 실험실}
\date{노트 457 · 2026년 8월 2일}
\maketitle

\begin{abstract}
깊이 $12$ 는 이 실험실에서 두 번 판정을 받았다. 노트 426 이 집 밖 짝
(앱 되뽑기 $1.000$ · KR $0.966$)으로 \textbf{채택}했고 --- 판을
$-0.0018$ 잃으면서 넣은 \emph{판정 지표를 거스른 첫 채택}이었다 ---
노트 432 가 \textbf{LODO $4/11$ · 평균 $-0.0046$} 으로 \textbf{되돌렸다}.

그 사이 노트 441 이 하나를 밝혔다. \textbf{LODO 열한 짝은 같은 수집 ·
같은 계기 · 같은 시기에서 나오므로 수집 단위로 눈금이 어긋나는 일을
원리상 못 본다.} 노트 456 이 K$=$64 로 같은 재개를 해 보고 조항을
만들었다 --- \emph{자가 바뀌면 그 자로 닫았던 갈래는 다시 열 수 있다.}

\textbf{되돌린 근거가 그 종류의 효과를 못 재는 자였다.} 다시 열었다.

\begin{center}
\begin{tabular}{lrr}
\hline
계기 & 깊이 $6 \to 12$ & 95\% 구간 \\
\hline
\textbf{판} (순위 지표) & $\mathbf{-0.0032}$ & $[-0.0041, -0.0023]$ \\
날짜 통제 판 & $-0.0046$ & --- \\
\textbf{KR 만화} (집 밖) & $\mathbf{+0.0064}$ & $[+0.0022, +0.0105]$ \\
\textbf{비게임 앱} (집 밖) & $\mathbf{+0.0061}$ & $[+0.0026, +0.0097]$ \\
\hline
\end{tabular}
\end{center}

씨앗 열을 \emph{짝지어} 쟀고 깊이 $8$ 이 네 계기 모두에서 중간이라
단조다. 노트 426 거울의 조건 셋이 다 찬다 --- ① 집 밖 짝 둘이 같은 부호
② \textbf{둘 다} 구간이 $0$ 을 안 뭄 ③ 판 손해 $0.0032$ 가 $2\times$짝SD
($0.0041$) 안. \textbf{채택한다.}
\end{abstract}

\begin{figure}[t]
\centering
\includegraphics[width=\textwidth]{figs/f1.pdf}
\caption{\textbf{가} --- 네 계기에서 깊이 $6 \to 12$. 판은 왼쪽, 집 밖
둘은 오른쪽. \textbf{나} --- 깊이 $8$ 이 전부 중간이라 우연한 봉우리가
아니다.}
\end{figure}

\section{거래를 숨기지 않는다}

\textbf{판은 확실히 잃는다.} 구간이 $0$ 을 안 물고, 짝 12뽑기에서
\textbf{양수 $1/12$} 이며, \textbf{날짜 통제 판도 같이 내려간다}
($-0.0046$) --- 노트 446 의 날짜 구멍 때문이 아니라 \emph{진짜 실력}을
잃는 것이다. 도메인으로는 웹툰 $-0.0086$($0/12$)과 팝업 $-0.0313$ 이
내고 모바일 $+0.0059$($12/12$)가 번다.

그것을 집 밖 두 짝의 $+0.006$ 씩과 바꾼다. 노트 444 의 단위로 옮기면
KR 은 천장의 $76.5\% \to 77.4\%$, 앱은 $41.8\% \to 42.8\%$ 다.

\textbf{이 거래가 옳은지는 규약이 정한다.} 노트 426 의 거울이 바로 이
상황을 위해 만들어졌고 --- 그때도 판 $-0.0018 \cdot 2/12$ 였다 ---
조건 셋이 다 찬다. 규약을 만들어 놓고 불편할 때만 안 따르면 규약이
아니다.

\section{왜 그때보다 다섯 배 작나}

노트 426 이 잰 전이 이득은 $+0.029$ 였고 지금은 $+0.006$ 이다.
\textbf{그 사이에 \texttt{OOC\_DROP} 이 들어갔다}(노트 441). 깊이가 사던
것의 상당 부분이 \emph{``눈금이 도메인마다 뒤집히는 축을 도메인별로
조건부로 쓰기''}였는데, 이제 그 축을 안 배운 수집에서는 \textbf{아예 안
쓴다}. 깊이가 하던 일을 더 싼 규칙이 가져갔고, 남은 $+0.006$ 이 깊이
고유의 몫이다.

\section{노트 428 의 어림이 여기서는 안 맞는다}

노트 428 은 \emph{관문(되뽑기)이 용량 변경을 아래로 어림한다}고 적었다
--- 되뽑기가 학습을 $63\%$ 로 줄이니 깊은 나무에 불리하다는 것이다.
그렇다면 관문의 판 손해가 전 자료보다 \emph{더} 나빠야 한다.

\textbf{반대다.} 관문 $-0.0028$, 전 자료 $-0.0032$. 크기는 작지만 방향이
어긋난다. 이번 한 번으로 노트 428 을 뒤집지는 않지만 \textbf{그 어림을
근거로 무엇을 봐주지도 않는다} --- 위 판정은 전 자료 수치($-0.0032$)로
했다.

\section*{되돌림 조항}

되돌릴 상태의 이름은 \textbf{``깊이 $6$ 인 F18''}(노트 432 규약).
되돌릴 조건도 미리 적는다.

\begin{quote}
\textbf{다음에 새로 여는 집 밖 짝에서 깊이 $12$ 의 이득이 $0$ 이하이거나,
판 손해가 다시 재서 $2\times$짝SD 밖으로 나가면 되돌린다.}
\end{quote}

새 챔피언 --- \textbf{판 $0.4828 \pm 0.0007$(SE · 씨앗 12)} · 날짜 통제 판
$0.5217$ · 분모 $3{,}430$ · 11 도메인.

\section*{남은 것}

\begin{itemize}
\item \textbf{깊이 $16 \cdot 24$ 는 안 봤다.} $8 \to 12$ 가 단조였으니
더 깊으면 더 버는지, 아니면 $12$ 가 봉우리인지.
\item \textbf{같은 잣대로 다시 열 것이 더 있나.} 노트 431 · 432 가 LODO
로 닫은 것 중 lr $0.10$ 은 판 근거로 남겼으니 해당 없고, 표지 둘째
(노트 420)는 \emph{전이를 잃어서} 닫은 것이라 자가 바뀌어도 방향이 같다.
\item \textbf{집 밖 짝 셋째} --- 짝이 둘뿐이라 ``둘 다 같은 부호''가
쉽게 찬다. 셋째 짝이 생기면 이 규약이 더 세진다.
\end{itemize}

\end{document}
